
% ==========================================
% ADDRESS MAPPING - K = 7126
% ==========================================

\clearpage
\begingroup

% Dinh dang rieng cho phan bang dia chi
\setstretch{1.05}
\setlength{\parindent}{0pt}
\setlength{\parskip}{0pt}
\renewcommand{\arraystretch}{1.3}
\setlength{\tabcolsep}{4pt}

% ==========================================
% PAGE 1
% ==========================================

\begin{center}
{\fontsize{15}{19}\selectfont\bfseries
BẢNG QUY HOẠCH ĐỊA CHỈ IP VÀ SỐ AS}
\end{center}

\vspace{0.35cm}

{\bfseries\large
1. Mã nhóm và quy tắc địa chỉ
}

\vspace{0.15cm}

Mã nhóm \textbf{K = 7126};
\textbf{X = 71}; \textbf{Y = 26}.
Các bảng dưới đây đã thay địa chỉ,
số AS và hostname theo mã nhóm.

\vspace{0.2cm}

\begin{table}[H]
\centering
\small
\begin{tabularx}{\textwidth}{|X|X|}
\hline
\textbf{Thành phần} & \textbf{Giá trị} \\
\hline
11 link giữa các router &
10.71.26.0/24, chia thành 11 mạng /30 \\
\hline
Loopback0 của router số n &
10.171.26.n/32 ($n = 1,\ldots,8$) \\
\hline
4 mạng LAN (Loopback1) &
10.171.126.0/24, chia thành 4 mạng /26 \\
\hline
Số AS của AS A, AS B, AS C &
64626, 64726, 64826 \\
\hline
Hostname &
\texttt{<tên router>-7126}
(ví dụ A1-7126, B2-7126) \\
\hline
\end{tabularx}
\end{table}

\vspace{0.2cm}

{\bfseries\large
2. Nối dây và địa chỉ
}

\vspace{0.15cm}

Mỗi router dùng tối đa 3 cổng:
Fa0/0, Fa0/1 (có sẵn) và Fa1/0
(card NM-1FE-TX ở slot 1).
Link Lk dùng mạng
10.71.26.[4 $\times$ (k -- 1)]/30;
router ghi trước lấy địa chỉ đầu,
router ghi sau lấy địa chỉ thứ hai.

\vspace{0.2cm}

\begin{table}[H]
\centering
\scriptsize
\renewcommand{\arraystretch}{1.5}
\setlength{\extrarowheight}{3pt}
\setlength{\tabcolsep}{3pt}

\begin{tabularx}{\textwidth}{
|c|l|l|l|>{\raggedright\arraybackslash}X|l|
}
\hline

\textbf{Link} &
\textbf{Router 1} &
\textbf{Router 2} &
\textbf{Mạng /30} &
\textbf{IP R1 / IP R2} &
\textbf{Loại} \\
\hline

L1 &
A1 (Fa0/0) &
A2 (Fa0/0) &
10.71.26.0 &
\makecell[l]{10.71.26.1\\10.71.26.2} &
RIPv2 \\
\hline

L2 &
A1 (Fa0/1) &
A3 (Fa0/0) &
10.71.26.4 &
\makecell[l]{10.71.26.5\\10.71.26.6} &
RIPv2 \\
\hline

L3 &
A2 (Fa0/1) &
A3 (Fa0/1) &
10.71.26.8 &
\makecell[l]{10.71.26.9\\10.71.26.10} &
RIPv2 \\
\hline

L4 &
A2 (Fa1/0) &
B1 (Fa0/0) &
10.71.26.12 &
\makecell[l]{10.71.26.13\\10.71.26.14} &
eBGP A--B \\
\hline

L5 &
A3 (Fa1/0) &
C1 (Fa0/1) &
10.71.26.16 &
\makecell[l]{10.71.26.17\\10.71.26.18} &
eBGP A--C \\
\hline

L6 &
B1 (Fa0/1) &
B2 (Fa0/0) &
10.71.26.20 &
\makecell[l]{10.71.26.21\\10.71.26.22} &
OSPF \\
\hline

L7 &
B2 (Fa0/1) &
B3 (Fa0/0) &
10.71.26.24 &
\makecell[l]{10.71.26.25\\10.71.26.26} &
OSPF \\
\hline

L8 &
B1 (Fa1/0) &
B3 (Fa0/1) &
10.71.26.28 &
\makecell[l]{10.71.26.29\\10.71.26.30} &
OSPF \\
\hline

L9 &
B3 (Fa1/0) &
C1 (Fa0/0) &
10.71.26.32 &
\makecell[l]{10.71.26.33\\10.71.26.34} &
eBGP B--C \\
\hline

L10 &
C1 (Fa1/0) &
D1 (Fa0/0) &
10.71.26.36 &
\makecell[l]{10.71.26.37\\10.71.26.38} &
\makecell[l]{Static\\Chính} \\
\hline

L11 &
B2 (Fa1/0) &
D1 (Fa0/1) &
10.71.26.40 &
\makecell[l]{10.71.26.41\\10.71.26.42} &
\makecell[l]{Static\\Dự phòng} \\
\hline

\end{tabularx}
\end{table}

{\footnotesize
\textit{
Tất cả 11 link dùng subnet mask
255.255.255.252 (/30).
Trong cột IP, dòng trên là Router 1,
dòng dưới là Router 2.
}
}

% ==========================================
% PAGE 2
% ==========================================

\clearpage

{\bfseries\large
3. Địa chỉ Loopback và LAN
}

\vspace{0.2cm}

\begin{table}[H]
\centering
\footnotesize
\renewcommand{\arraystretch}{1.65}
\setlength{\tabcolsep}{4pt}

\begin{tabularx}{\textwidth}{
|l|c|>{\raggedright\arraybackslash}X|l|
>{\raggedright\arraybackslash}p{4.1cm}|
}
\hline

\textbf{Router} &
\textbf{n} &
\textbf{Vai trò} &
\textbf{Loopback0} &
\textbf{Loopback1 = LAN} \\
\hline

A1-7126 & 1 &
AS A, nội bộ &
10.171.26.1 &
LAN-A: 10.171.126.1
\newline Mạng 10.171.126.0/26 \\
\hline

A2-7126 & 2 &
AS A, biên, nối ISP B &
10.171.26.2 &
-- \\
\hline

A3-7126 & 3 &
AS A, biên, nối ISP C &
10.171.26.3 &
-- \\
\hline

B1-7126 & 4 &
AS B, biên, nối AS A &
10.171.26.4 &
-- \\
\hline

B2-7126 & 5 &
AS B, nội bộ, nối khách hàng D1 &
10.171.26.5 &
LAN-B: 10.171.126.65
\newline Mạng 10.171.126.64/26 \\
\hline

B3-7126 & 6 &
AS B, biên, nối AS C &
10.171.26.6 &
-- \\
\hline

C1-7126 & 7 &
AS C (ISP), biên &
10.171.26.7 &
LAN-C: 10.171.126.129
\newline Mạng 10.171.126.128/26 \\
\hline

D1-7126 & 8 &
Khách hàng, không chạy giao thức định tuyến &
10.171.26.8 &
LAN-D: 10.171.126.193
\newline Mạng 10.171.126.192/26 \\
\hline

\end{tabularx}
\end{table}

{\footnotesize
\textit{
Loopback0 dùng subnet mask
255.255.255.255 (/32);
Loopback1 dùng 255.255.255.192 (/26).
Router không có LAN không cấu hình Loopback1.
}
}

\vspace{0.4cm}

{\bfseries\large
4. Tóm tắt bốn mạng LAN
}

\vspace{0.2cm}

\begin{table}[H]
\centering
\small
\begin{tabularx}{\textwidth}{
|c|l|X|X|
}
\hline

\textbf{LAN} &
\textbf{Router} &
\textbf{Mạng LAN (/26)} &
\textbf{IP Loopback1 (/26)} \\
\hline

LAN-A & A1-7126 &
10.171.126.0/26 &
10.171.126.1/26 \\
\hline

LAN-B & B2-7126 &
10.171.126.64/26 &
10.171.126.65/26 \\
\hline

LAN-C & C1-7126 &
10.171.126.128/26 &
10.171.126.129/26 \\
\hline

LAN-D & D1-7126 &
10.171.126.192/26 &
10.171.126.193/26 \\
\hline

\end{tabularx}
\end{table}

\vspace{0.4cm}

{\bfseries\large
5. Địa chỉ peer eBGP và next-hop static
}

\vspace{0.2cm}

\begin{table}[H]
\centering
\small
\begin{tabularx}{\textwidth}{
|l|X|X|
}
\hline

\textbf{Link} &
\textbf{Phía 1} &
\textbf{Phía 2} \\
\hline

L4 -- eBGP A--B &
A2-7126: 10.71.26.13 &
B1-7126: 10.71.26.14 \\
\hline

L5 -- eBGP A--C &
A3-7126: 10.71.26.17 &
C1-7126: 10.71.26.18 \\
\hline

L9 -- eBGP B--C &
B3-7126: 10.71.26.33 &
C1-7126: 10.71.26.34 \\
\hline

L10 -- Static chính &
C1-7126: 10.71.26.37 &
D1-7126: 10.71.26.38 \\
\hline

L11 -- Static dự phòng &
B2-7126: 10.71.26.41 &
D1-7126: 10.71.26.42 \\
\hline

\end{tabularx}
\end{table}

\vspace{0.2cm}

\vfill


\endgroup
\clearpage
